\chapter{Software realizzato}
\label{app:codice}

Questa appendice documenta l'organizzazione del software prodotto per la tesi e i
comandi con cui è stato usato. I listati completi non sono riportati integralmente:
si indicano struttura, scelte implementative rilevanti e modalità d'uso. Tutto il
software è nel repository della tesi insieme ai file CSV della campagna, che ne sono
l'unico contenuto non ricostruibile.

\section{Organizzazione del repository}

\begin{center}\small
\begin{tabularx}{\linewidth}{@{}lX@{}}
\toprule
\textbf{Percorso} & \textbf{Contenuto} \\
\midrule
\multicolumn{2}{@{}l}{\emph{Firmware per ESP32 (Arduino)}} \\
\texttt{firmware/ld2410b\_logger/} & logger CSV a 5~Hz con engineering mode e
  lettura del PIR: il firmware delle fasi 1--7 sul \ldb \\
\texttt{firmware/ld2410b\_web/} & la web UI del \cref{cap:webui}: Access Point,
  server asincrono, WebSocket, indice di vitalità a bordo; stesso CSV del logger
  sulla seriale, usato per la Fase~8 \\
\texttt{firmware/ld2420\_logger\_bin/} & logger per il \ldc\ in modalità binaria:
  presenza, distanza e 16 energie a 16~bit a 10~Hz, con scrittura opzionale di gate
  massimo, gate minimo e ritardo a ogni avvio \\
\texttt{firmware/test00\_ld2410\_info/} & lettura dei parametri del \ldb\ via UART:
  firmware, MAC, soglie, gate, timeout \\
\texttt{firmware/test01\_ld2410\_base/} & sketch diagnostico con auto-scan del baud
  rate e stampa dei byte grezzi \\
\texttt{firmware/test03\_pir\_base/} & verifica del PIR e misura della durata degli
  impulsi \\
\texttt{firmware/test04\_set\_gate/} & configurazione del gate massimo del \ldb\ per
  la prova di selettività spaziale \\
\texttt{firmware/test06\_ld2420\_set\_gate/} & lettura e scrittura dei parametri del
  \ldc\ via UART (32 soglie, gate, ritardo) \\
\texttt{firmware/ld2420\_monitor/}, \texttt{test05\_ld2420\_raw/} & lettura del \ldc\
  in modalità ASCII e dump dei byte grezzi (superati dal logger binario) \\
\texttt{tools/embed\_web.py} & comprime la pagina web in gzip e genera l'intestazione
  C con i byte in flash \\
\midrule
\multicolumn{2}{@{}l}{\emph{Acquisizione}} \\
\texttt{HLK-LD2410x/acquire.py} & acquisizione dalla seriale con metadati
  sperimentali (versione adattata dal repository di riferimento) \\
\texttt{HLK-LD2410x/data/} & i CSV della campagna e il registro delle sessioni \\
\texttt{HLK-LD2420/Backup config/} & le configurazioni del \ldc\ salvate dallo
  strumento del produttore (fabbrica e tarate), in XML \\
\bottomrule
\end{tabularx}
\end{center}

\begin{center}\small
\begin{tabularx}{\linewidth}{@{}lX@{}}
\toprule
\textbf{Percorso} & \textbf{Contenuto} \\
\midrule
\multicolumn{2}{@{}l}{\emph{Analisi (Python)}} \\
\texttt{analisi/verifica\_engineering.py} & controllo di qualità dei file CSV \\
\texttt{analisi/analizza\_test.py} & metriche di confronto e aggregazione per
  scenario \\
\texttt{analisi/analizza\_respiro.py} & analisi in frequenza e ricerca del canale
  migliore \\
\texttt{analisi/portata2420.py} & analisi per finestre temporali delle prove di
  portata del \ldc \\
\texttt{analisi/vitalita\_proto.py} & prototipo dell'indice di vitalità (v2 e v3),
  taratura e matrice di confusione \\
\texttt{analisi/verifica\_vitalita\_bordo.py} & confronto fra l'indice calcolato a
  bordo e il ricalcolo offline \\
\midrule
\multicolumn{2}{@{}l}{\emph{Materiali derivati}} \\
\texttt{analisi/rigenera\_tutto.py} & un comando per rifare figure, foglio Excel,
  pagina di riepilogo e PDF dai CSV \\
\texttt{analisi/grafici\_tesi.py}, \texttt{grafici\_tesi\_2.py} & le 22 figure di
  questa tesi, in PDF vettoriale e PNG \\
\texttt{analisi/esporta\_excel.py} & foglio Excel con una riga per trial valido e i
  dati aggregati per figura \\
\texttt{analisi/genera\_pagina.py} & pagina HTML di riepilogo autonoma, con le
  figure incorporate \\
\bottomrule
\end{tabularx}
\end{center}

\section{Il logger del \ldb}

\subsection{Scelte implementative rilevanti}

\begin{description}[leftmargin=1.4em,style=nextline,itemsep=3pt]
  \item[Periodo fisso e non ritardo fisso] Il ciclo principale calcola l'istante del
    prossimo campione anziché attendere un intervallo dopo l'elaborazione. È la
    ragione per cui la cadenza misurata è di 200~ms esatti su tutti gli intervalli,
    senza deriva: requisito necessario per l'analisi in frequenza.
  \item[Lettura del PIR nello stesso frame] Il valore del PIR è campionato
    immediatamente prima di comporre la riga CSV, quindi i due sensori condividono la
    base dei tempi per costruzione.
  \item[Engineering mode all'avvio] Il firmware entra in configurazione, attiva
    l'engineering mode, esce dalla configurazione e verifica che i frame ricevuti
    siano del tipo esteso, segnalando l'esito sulla seriale. Un fallimento silenzioso
    produrrebbe file senza colonne per gate scoperti solo in fase di analisi.
  \item[Intestazione stampata una sola volta] L'intestazione CSV è stampata in
    \texttt{setup()}. Poiché il reset all'apertura della porta non è garantito, lo
    script di acquisizione la ricostruisce dal numero di colonne (si veda oltre).
  \item[Marcatore di versione] Ogni sketch stampa all'avvio una riga
    \texttt{\# BUILD n}, che \texttt{acquire.py} scarta perché inizia con il cancelletto
    ma che resta nel log: l'esito di una prova non può così essere attribuito a un
    binario diverso da quello caricato.
\end{description}

\subsection{Formato di uscita}

\begin{lstlisting}[caption={Intestazione CSV completa in engineering mode, 29
colonne.},basicstyle=\ttfamily\scriptsize]
timestamp_ms,radar_presence,moving_target,stationary_target,
moving_distance_cm,stationary_distance_cm,moving_energy,stationary_energy,
pir_presence,
menergy_gate0,menergy_gate1,menergy_gate2,menergy_gate3,menergy_gate4,
menergy_gate5,menergy_gate6,menergy_gate7,menergy_gate8,
senergy_gate0,senergy_gate1,senergy_gate2,senergy_gate3,senergy_gate4,
senergy_gate5,senergy_gate6,senergy_gate7,senergy_gate8,
light_level,out_level
\end{lstlisting}

Le prime nove colonne coincidono con quelle del repository di
riferimento~\cite{repoCallisto}, per compatibilità. Il firmware della web UI produce
sulla seriale esattamente queste 29 colonne, quindi \texttt{acquire.py} funziona con
entrambi senza modifiche.

\section{Il logger del \ldc}

Il logger binario legge i frame da 45~byte della modalità \emph{energy}
(\cref{app:ld2420-frame}) a 10~Hz e li riduce a 5~Hz per coerenza con il resto della
campagna. Produce le stesse nove colonne di base, seguite dalle sedici energie grezze
(\texttt{energy2420\_gate0..15}, interi 0--65535: il nome è diverso da quello del
\ldb\ proprio perché la scala è diversa e nessuno script le confonda), dal conteggio
dei frame validi e dalla distanza grezza. Le colonne del \ldb\ che il \ldc\ non
possiede (\texttt{moving\_energy}, \texttt{stationary\_*}) valgono zero.

Due scelte meritano una nota.

\begin{description}[leftmargin=1.4em,style=nextline,itemsep=3pt]
  \item[Il PIR non collegato vale $-1$, non $0$] Nei test del solo \ldc\ il PIR non è
    cablato. Un piedino non collegato leggerebbe zero per tutto il file, e gli script
    ne ricaverebbero un ``$100\,\%$ di falsi negativi'' con una persona davanti: un
    numero finto che sembra una misura. Il firmware scrive $-1$, lo script di
    analisi omette tutte le metriche del PIR e marca lo stato come non collegato, e
    nel foglio Excel le celle restano vuote con la ragione a fianco. Verificato che
    sui file con PIR cablato nulla cambi.
  \item[Parametri scritti a ogni avvio, in RAM] Gate massimo, gate minimo e ritardo
    di scomparsa possono essere scritti via UART a ogni accensione dell'ESP32. La
    scrittura è accettata e riletta correttamente, ma non è persistente: il modulo la
    perde al proprio ciclo di alimentazione. Il valore riletto finisce in un
    commento in testa al CSV, così ogni file dichiara la configurazione con cui è
    stato acquisito.
\end{description}

\section{Acquisizione}

\subsection{Ambiente}

\begin{lstlisting}[language=bash,caption={Preparazione dell'ambiente Python.},
basicstyle=\ttfamily\scriptsize]
cd HLK-LD2410x
python -m venv .venv
.venv\Scripts\Activate.ps1
pip install pyserial numpy
mkdir data
\end{lstlisting}

Gli script dei materiali derivati richiedono inoltre \texttt{matplotlib},
\texttt{openpyxl} e \texttt{pillow}; \texttt{rigenera\_tutto.py} lo verifica e indica
il comando di installazione se mancano.

\subsection{Uso}

\begin{lstlisting}[language=bash,caption={Esempio di acquisizione di un trial.},
basicstyle=\ttfamily\scriptsize]
python acquire.py --port COM3 --duration 300 ^
  --output data/fermo_seduto_T01.csv ^
  --scenario fermo_seduto --trial T01 ^
  --ground_truth_presence 1 --ground_truth_state still
\end{lstlisting}

Lo script aggiunge a ogni riga le colonne \texttt{pc\_time\_s},
\texttt{group\_id}, \texttt{trial\_id}, \texttt{scenario},
\texttt{ground\_truth\_presence}, \texttt{ground\_truth\_state}. Le opzioni
\texttt{-{}-beep-at SEC} (segnale acustico all'istante dell'evento, per le latenze)
e \texttt{-{}-tappe} (annunci vocali a istanti prefissati, per le prove di portata)
sono aggiunte per questa tesi: il segnale conta dalla prima riga di dati, cioè dalla
stessa origine dei tempi usata dall'analisi, mentre un timer esterno introdurrebbe
un errore sistematico di 2--3~s pari al reset dell'ESP32.

\subsection{La modifica necessaria}

La versione originale scrive righe solo dopo aver riconosciuto l'intestazione
stampata dall'ESP32 in \texttt{setup()}. Se il reset all'apertura della porta non
scatta, il file risulta \textbf{vuoto senza alcun messaggio}. La versione adattata
ricostruisce l'intestazione dal numero di colonne della prima riga di dati:

\begin{center}\small
\begin{tabular}{@{}rp{0.72\linewidth}@{}}
\toprule
\textbf{Colonne} & \textbf{Interpretazione} \\
\midrule
9 & modalità base \\
27 & engineering mode del \ldb\ (o logger binario del \ldc, che stampa l'intestazione a ogni avvio) \\
29 & engineering mode con sensore di luce e stato OUT \\
\bottomrule
\end{tabular}
\end{center}

\section{Script di analisi}

\subsection{Controllo di qualità}

\begin{lstlisting}[language=bash,basicstyle=\ttfamily\scriptsize]
python analisi/verifica_engineering.py data/fermo_seduto_T01.csv
\end{lstlisting}

Verifica ed esce con diagnostica su: cadenza effettiva e jitter, presenza delle
colonne per gate, percentuale di campioni saturi per canale, coerenza fra gate di
picco e distanza riportata, rumore di fondo per gate nei tratti a stanza vuota,
numero e istante delle transizioni di presenza. Va eseguito su ogni file
\emph{prima} di usarlo per l'analisi.

\subsection{Metriche di confronto}

\begin{lstlisting}[language=bash,basicstyle=\ttfamily\scriptsize]
python analisi/analizza_test.py data/fermo_seduto_T0*.csv --salta-inizio 20
python analisi/analizza_test.py data/ingresso_T0*.csv --event-time 30
python analisi/analizza_test.py data/uscita_T0*.csv --release-time 30
\end{lstlisting}

Le opzioni principali: \texttt{-{}-salta-inizio} scarta i primi secondi (uscita
dell'operatore e stabilizzazione del PIR; 20~s negli scenari ordinari, 40 sotto il
banco, 120 nella selettività, 60 a stanza vuota, 90 nei trial da fermo del \ldc),
\texttt{-{}-event-time} e \texttt{-{}-release-time} attivano il calcolo delle latenze
con la convenzione dell'evento a istante noto. Con più file dello stesso scenario lo
script aggrega e riporta media e deviazione standard fra trial. Nel calcolo della
latenza di rilascio distingue i casi in cui il rilascio non avviene mai e quelli in
cui avviene prima dell'istante dichiarato, e conta le riaccensioni nella coda. Lo
script legge indifferentemente i file del \ldb, del \ldc\ e della web UI.

\subsection{Analisi del respiro}

\begin{lstlisting}[language=bash,basicstyle=\ttfamily\scriptsize]
python analisi/analizza_respiro.py data/respiro_2m_15_T01.csv --scan ^
  --salta-inizio 40 --min-canali 3
\end{lstlisting}

La modalità \texttt{-{}-scan} esamina tutti i canali di energia, scarta quelli saturi
e quelli piatti secondo criteri dichiarati, ordina i rimanenti per rapporto
segnale-rumore del picco in banda 0,1--0,5~Hz, raggruppa le stime concordi entro il
$10\,\%$ e applica le regole di \cref{sec:respiro-metronomo}: almeno $N$ canali in
movimento nel gruppo, preferenza per la fondamentale fra gruppi in rapporto armonico,
soglia di plausibilità a 8 atti al minuto. Esporta lo spettro in CSV. Richiede NumPy;
gli altri script di analisi usano la sola libreria standard.

\subsection{Indice di vitalità}

\begin{lstlisting}[language=bash,basicstyle=\ttfamily\scriptsize]
python analisi/vitalita_proto.py data/fermo_1m_H_T0*.csv ^
  data/micromovimenti_1m_H_T0*.csv data/movimento_1m_H_T0*.csv ^
  --alpha-mov 0.05 --alpha-var 0.01 --k 0.5 --soglie 45,95 ^
  --gate distanza --fondo-da data/stanza_vuota_notte_T01.csv
python analisi/verifica_vitalita_bordo.py data/Immobile_web_1m_T01.csv
\end{lstlisting}

Il prototipo calcola la serie dell'indice per ogni file, la distribuzione per
scenario e la matrice di confusione scenario/classe; le opzioni \texttt{-{}-gate}
(\texttt{energia}, \texttt{distanza} o un gate fisso) e \texttt{-{}-sorgente}
(\texttt{gate} o \texttt{max}) permettono di riprodurre le versioni precedenti
dell'algoritmo per confronto. Lo script di verifica rilegge un CSV esportato dalla
web UI, ricalcola l'indice con le stesse costanti e lo confronta con la colonna
scritta dal firmware, riportando la frazione di campioni concordi entro $\pm 1$
sull'intero file e sulla coda a regime.

\section{Materiali derivati riproducibili}

\begin{lstlisting}[language=bash,basicstyle=\ttfamily\scriptsize]
python analisi/rigenera_tutto.py
\end{lstlisting}

Esegue in ordine i due script delle figure, l'esportazione Excel, la pagina di
riepilogo e la stampa in PDF, in circa mezzo minuto, fermandosi al primo errore senza
sovrascrivere i passi successivi. Le figure importano le funzioni di
\texttt{analizza\_test.py} anziché ricalcolare: i numeri nei grafici coincidono per
costruzione con quelli dei capitoli. Figure, foglio Excel, pagina e PDF non sono
versionati perché si rifanno dai CSV; il foglio Excel ha una riga per trial valido e
un foglio separato con i file esclusi e la ragione dell'esclusione, così che nessun
file di prova o di controllo entri nelle statistiche per una svista.

\section{Configurazione dei parametri via UART}

Lo sketch \texttt{firmware/test04\_set\_gate/} imposta il gate massimo del \ldb\ per
la prova di selettività spaziale. Tre accorgimenti sono degni di nota, perché
riguardano la scrittura di parametri persistenti su un modulo:

\begin{enumerate}[leftmargin=1.8em,itemsep=2pt]
  \item dopo ogni scrittura \textbf{rilegge sempre} i parametri per conferma, e
        segnala se il valore letto non corrisponde a quello scritto;
  \item \textbf{avvisa} quando la configurazione corrente è diversa da quella di
        fabbrica, per evitare di condurre acquisizioni con parametri modificati
        inconsapevolmente --- il rischio più insidioso in una campagna comparativa;
  \item il comando \texttt{d} riscrive i valori di fabbrica del nostro esemplare
        (le 9+9 soglie di \cref{tab:identita}, gate 8/8 e timeout di 5~s),
        così da poter tornare in ogni momento alla configurazione di riferimento.
\end{enumerate}

Lo sketch \texttt{firmware/test06\_ld2420\_set\_gate/} fa lo stesso per il \ldc:
legge le 32 soglie, il gate massimo e il ritardo, e ha permesso di verificare che i
valori in flash coincidessero con quelli mostrati in decibel dallo strumento del
produttore, secondo la relazione $\text{grezzo} = 10^{\text{dB}/10}$
(\cref{sec:ld2420-stato}). Per il \ldc\ la scrittura via UART non è persistente, e
la configurazione di riferimento resta quella salvata in flash dallo strumento del
produttore, di cui i file XML nel repository sono la copia.

\section{Come si compila la web UI}

\begin{enumerate}[leftmargin=1.8em,itemsep=2pt]
  \item installare dal gestore di librerie dell'IDE \texttt{ESP Async WebServer} e
        \texttt{Async TCP} nel fork ESP32Async, oltre a \texttt{MyLD2410} e
        \texttt{ArduinoJson}~7;
  \item in \emph{Strumenti $\to$ Partition Scheme} scegliere ``Huge APP (3MB No
        OTA)'';
  \item dopo ogni modifica ai file in \texttt{web/}, eseguire
        \texttt{python tools/embed\_web.py}, che rigenera \texttt{web\_assets.h};
  \item compilare e caricare come uno sketch qualsiasi. All'avvio la seriale a
        115200~baud stampa il marcatore di versione, poi le righe CSV.
\end{enumerate}

La rete creata è \texttt{DIPME-Sensor}; la dashboard risponde a
\url{http://192.168.4.1/} e a qualunque nome, grazie al DNS a bordo. Per
scaricare il CSV va usato un browser vero e non la finestra del portale captive.
